% -----------------------------------------------
% ISMIR 2026 LBD — Pocket Producer
% Based on the official ISMIR2026_lbd_template.tex
% (sha256 e368d536b41c...dec3d, fetched 2026-07-31).
% Rules: 2 pages scientific content + 1 page references only.
% LBD submissions are NOT anonymized.
% -----------------------------------------------

\documentclass{article}
\usepackage[T1]{fontenc}
\usepackage[utf8]{inputenc}
\usepackage[lbd,submission]{ismir}
\usepackage{amsmath,cite,url}
\usepackage{microtype}
\usepackage{graphicx}
\usepackage{color}
\usepackage{booktabs}
\usepackage{tikz}
\usetikzlibrary{arrows.meta,positioning,decorations.pathreplacing}
\def\lbd{}

\title{Pocket Producer: Session-Conditioned Retrieval\\
as User-Controlled Proposal}

\oneauthor
  {Chuling Li} {University of Nottingham \\ {\tt scycl14@nottingham.ac.uk}}

\def\authorname{C. Li}

\sloppy

\begin{document}

\maketitle

\begin{abstract}
We present Pocket Producer, a working prototype that operationalizes
\emph{session-conditioned continuation retrieval} inside the Audiotool
DAW\@.  A five-signal logistic fusion, trained with BPR on pack
co-membership weak labels from the Freesound Loop Dataset, is deployed
in the system as a cold-start candidate generator.
Four lines of evidence show that this proxy should not be trusted to
decide autonomously:
(1)~the apparent fusion gain is label-sensitive---+4.0\,pp over frozen
cosine on hard\_similar negatives when 26\% of weak labels group items
by uploader rather than pack, but $-0.03$\,pp under pack-only labels;
(2)~leave-one-signal-out ablation reveals regime-dependent signal
contributions---tempo helps on hard\_similar but hurts on
hard\_tempo\_key;
(3)~on a held-out split committed to version control before
evaluation, the fusion$-$cosine difference is small and uncertain
(hard\_similar $+1.5$\,pp, 95\%~CI crossing zero) and significantly
negative on hard\_tempo\_key;
(4)~an exploratory pilot shows a descriptive reversal between
weak-label signal rankings and human consensus.
Pocket Producer therefore treats every recommendation as a
\emph{proposal}: candidates are previewable, insertion is reversible,
and feedback is logged for future personalization---which is not yet
evaluated.
\end{abstract}

% ==================================================================
\section{Introduction}\label{sec:intro}
% ==================================================================

When a musician returns to an unfinished loop-based project, the
question is not ``what sounds similar?''\ but ``what would be useful
as the next step?''
Loop compatibility~\cite{chen2020neuralloop,lattner2022samplematch,yi2024drum2bass,roychaudhuri2026loopmatcher},
stem coherence~\cite{ciranni2025cocola,nistal2024diffariff,riou2025stemretrieval},
session-conditioned generation~\cite{he2025tomi}, and reversible
co-creation~\cite{elkins2025dawzy} all address parts of this problem.
We define \textbf{continuation utility} as the property that a
candidate helps the musician make progress in the current
session---unlike similarity, it is asymmetric, context-dependent, and
multi-dimensional.

We contribute: (1)~a working DAW-integrated prototype whose deployed
ranker is the same fusion model evaluated here, treating retrieval as
user-controlled proposal, not automatic decision;
(2)~empirical evidence that weak-label proxy scores are sensitive to
label construction, regime-dependent, uncertain under held-out
evaluation, and may not reflect human judgment---motivating the
design choice to keep creative authority with the musician.

% ==================================================================
\section{System}\label{sec:system}
% ==================================================================

\begin{figure}[t]
  \centering
  \resizebox{\columnwidth}{!}{% TikZ system pipeline — Pocket Producer (revised: user-control emphasis)
\definecolor{csession}{HTML}{9BCAE7}
\definecolor{cfingerprint}{HTML}{D4E6A9}
\definecolor{cfusion}{HTML}{D4BFDC}
\definecolor{cpreview}{HTML}{F9D8A0}
\definecolor{caction}{HTML}{F9D1D2}
\definecolor{cbadge}{HTML}{477DC0}
\definecolor{cfeedback}{HTML}{477DC0}
\definecolor{cedge}{HTML}{3A3A3A}
\definecolor{ctext}{HTML}{1A1A1A}
\definecolor{csub}{HTML}{555555}

\begin{tikzpicture}[
    node distance=0.4cm and 0.45cm,
    every node/.style={font=\footnotesize},
    sysbox/.style 2 args={
        rectangle, rounded corners=2.5pt,
        minimum width=2.1cm, minimum height=0.85cm,
        draw=cedge, line width=0.5pt,
        fill=#1, text=ctext, align=center,
        inner sep=3pt,
        append after command={
            node [circle, fill=cbadge, text=white,
                  font=\sffamily\tiny\bfseries,
                  inner sep=0pt, minimum size=10pt]
            at ([shift={(-0.10cm,0.10cm)}]\tikzlastnode.north west) {#2}
        },
    },
    sub/.style={font=\tiny, text=csub},
    arr/.style={
        -{Stealth[length=4pt, width=3.5pt]},
        line width=0.55pt, color=cedge,
    },
    fbarr/.style={
        -{Stealth[length=4pt, width=3.5pt]},
        line width=0.55pt, color=cfeedback, dashed,
    },
    brace/.style={
        decorate, decoration={brace, amplitude=4pt, mirror},
        line width=0.4pt, color=csub,
    },
]

% Row 1 (top): candidate generation
\node[sysbox={csession}{1}]   (S) {\textbf{DAW Session}};
\node[sysbox={cfingerprint}{2}, right=of S] (F) {\textbf{Fingerprint}};
\node[sysbox={cfusion}{3}, right=of F]   (Fu) {\textbf{Multi-Signal}\\\textbf{Ranking}};

% Row 2 (bottom): user control
\node[sysbox={cpreview}{4}, below=0.9cm of Fu]  (P) {\textbf{Preview}};
\node[sysbox={caction}{5}, left=of P]       (A)  {\textbf{Insert / Undo}};
\node[sysbox={csession!30}{6}, left=of A]         (Fb)  {\textbf{Feedback Log}};

% Sublabels
\node[sub, below=1pt of S.south]  {Audiotool Nexus};
\node[sub, below=1pt of F.south]  {tempo, key, embeddings};
\node[sub, below=1pt of Fu.south]  {cold-start BPR fusion};
\node[sub, above=1pt of P.north] {audition before commit};
\node[sub, above=1pt of A.north]  {single-transaction};
\node[sub, above=1pt of Fb.north]  {accept / reject / undo};

% Forward arrows
\draw[arr] (S.east)  -- (F.west);
\draw[arr] (F.east)  -- (Fu.west);
\draw[arr] (Fu.south) -- (P.north);
\draw[arr] (P.west) -- (A.east);
\draw[arr] (A.west) -- (Fb.east);

% Feedback loop
\coordinate (fb) at ([xshift=-0.5cm]Fb.west);
\draw[fbarr] (Fb.west) -- (fb) |- (S.west);
\node[font=\tiny, text=cfeedback, rotate=90, anchor=south]
    at ([xshift=-0.22cm]fb) {\textit{future personalization}};

% Brace: proposer vs decider
\draw[brace] ([yshift=-0.35cm]S.south west) -- ([yshift=-0.35cm]Fu.south east)
    node[midway, below=5pt, font=\tiny\itshape, text=csub] {system proposes};
\draw[brace] ([yshift=0.35cm]Fb.north west) -- ([yshift=0.35cm]P.north east)
    node[midway, above=5pt, font=\tiny\itshape, text=csub] {musician decides};

\end{tikzpicture}
}
  \caption{Pocket Producer pipeline: the system proposes
  (steps~1--3); the musician decides (steps~4--6).  Feedback is
  logged but personalization is not yet trained.}
  \label{fig:system}
\end{figure}

Figure~\ref{fig:system} shows the pipeline.
The live Audiotool session is read via the Nexus SDK and mapped to
a session fingerprint (tempo, active regions, CLAP
embeddings~\cite{elizalde2023clap}).
A five-signal logistic fusion trained with BPR
loss~\cite{rendle2009bpr} ranks candidates across timbral similarity
(CLAP cosine to session regions, mean- and max-pooled), tempo match,
key compatibility, and tag overlap.
The deployed ranker is this same fusion model, registered as the
default whenever session region audio is available, with two stated
deltas from the offline evaluation: Audiotool documents carry no key
signature, so the key term is constant zero online, and tag context
comes from fragments used in the session rather than pack siblings.
Without region audio the service falls back to a transparent rules
baseline and reports that it did so.
The system positions itself as a \emph{proposer}: candidates are
surfaced for preview before any session state changes.
Accepted candidates are inserted via a single Nexus transaction with
undo support; rejected or undone events are logged.
This design is a response to the empirical findings in
Section~\ref{sec:evidence}: because the proxy's correspondence to
continuation utility is limited and regime-dependent, the system must
not substitute its ranking for the musician's judgment.

% ==================================================================
\section{Weak-Label Fusion}\label{sec:fusion}
% ==================================================================

From the FSLD~\cite{ramires2020fsld} we kept only sounds with a true
Freesound pack (1,797 items, 605 packs), discarding items whose only
grouping is a shared uploader.
Each multi-item pack yields one leave-one-out query per item, with
same-pack siblings as positives and cross-pack items as negatives in
three difficulty tiers (\emph{easy}: tempo/key mismatch;
\emph{hard\_tempo\_key}: compatible tempo and key;
\emph{hard\_similar}: highest cosine to the positive from another
pack).
A 121-source held-out split (seed 20260810) was committed to version
control together with the evaluation protocol \emph{before} any
evaluation ran; all development below uses only the 484 train
sources (981 examples).
LAION-CLAP Music~\cite{wu2023laionclap} (512-dim) provides the
embedding backbone.

On the train side (mean over 5 source-grouped splits), frozen cosine
reaches 0.929 overall and 0.814 on hard\_similar; five-signal fusion
reaches 0.925 and 0.813.
Under pack-only labels the fusion advantage is gone: $-0.03$\,pp on
hard\_similar, and on hard\_tempo\_key fusion is significantly
\emph{worse} than cosine ($-2.1$\,pp, paired bootstrap 95\%~CI
$[-3.9, -0.6]$\,pp).

% ==================================================================
\section{Why the Proxy Should Not Decide Alone}%
\label{sec:evidence}
% ==================================================================

Four lines of evidence, each bearing a limited conclusion, jointly
motivate treating the proxy as a proposal rather than a decision.

\subsection{The Gain Is Label-Sensitive}\label{sec:labels}

An earlier version of this evaluation used the manifest's fallback
grouping, under which 26\% of examples treated
\emph{same-uploader} items as co-members of a production context.
On that mixed-label corpus, fusion appeared to beat frozen cosine by
+4.0\,pp on hard\_similar (949 sources, same protocol).
Restricting to true pack co-membership erases the gain
(Section~\ref{sec:fusion}).
A proxy whose measured advantage depends this strongly on how weak
labels are constructed cannot be treated as a reliable estimate of
continuation utility.

\subsection{Regime-Dependent Signal Contributions}\label{sec:loso}

\begin{table}[t]
  \centering
  \caption{LOSO ablation: accuracy change (pp) when one signal is
  removed.  Pack-only train side, 484 sources, mean over 5 splits.
  Tempo (bold) shows the clearest regime dependence.}
  \label{tab:loso}
  \footnotesize
  % LOSO delta table (pp): removing signal → accuracy change
% Pack-only development set (484 train sources), mean over 5 source-grouped splits
% Negative values = removing hurts (signal contributes)
\setlength{\tabcolsep}{4pt}
\begin{tabular}{@{}l r r r@{}}
\toprule
\textbf{Signal removed} & \textbf{easy} & \textbf{hard\_sim.} & \textbf{hard\_tk} \\
\midrule
audio\_cos\_mean & $-0.32$ & $-0.36$ & $-1.04$ \\
audio\_cos\_max  & $+0.17$ & $-0.30$ & $-0.15$ \\
\textbf{tempo}   & $-0.27$ & $\mathbf{-0.96}$ & $\mathbf{+0.68}$ \\
key              & $-0.21$ & $+0.94$ & $+1.21$ \\
tag\_jaccard     & $+0.09$ & $-0.79$ & $-0.07$ \\
\bottomrule
\end{tabular}

\end{table}

Table~\ref{tab:loso} shows leave-one-signal-out deltas per negative
regime.  Removing tempo costs $-0.96$\,pp on hard\_similar (where
candidates sound alike and tempo becomes a distinguishing factor)
but \emph{gains} $+0.68$\,pp on hard\_tempo\_key; removing key helps
on both hard regimes ($+0.94$\,/\,$+1.21$\,pp).
Both patterns replicate directionally on the held-out split (tempo
$-1.48$\,/\,$+0.37$\,pp, key $+0.62$\,/\,$+0.97$\,pp, retrained on
the train side).
Part of this pattern is built in: the easy and hard\_tempo\_key
regimes are themselves constructed from tempo and key, so a
tempo/key signal is near-constant within them by design.
The point is not that the dependence is surprising, but that a
single fused score averages over regimes in which individual signals
switch sign---so its ranking cannot be uniformly trusted.

\subsection{Held-Out Evaluation}\label{sec:heldout}

The 121-source held-out split was evaluated exactly once, after all
development; protocol and split were committed to version control
first, so the freeze is verifiable in the repository history.
Fusion was retrained on the 484-source train side with the same
hyperparameters and 5 seeds.
The fusion$-$cosine difference on hard\_similar is $+1.5$\,pp with a
source-grouped bootstrap 95\%~CI of $[-1.5, +4.7]$\,pp (joint
resampling shared across seeds)---small and uncertain.
On hard\_tempo\_key, fusion performs significantly \emph{worse} than
cosine ($-1.0$\,pp, CI entirely below zero); overall the difference
is $+0.4$\,pp with a CI crossing zero.

We note residual limitations: the signal set and hyperparameters were
designed on the earlier mixed-label corpus, and the pack-only corpus
is a subset of data inspected during that development.  This run is a
sensitivity re-analysis under a committed split, not de novo
confirmatory evidence.

\subsection{Exploratory Pilot}\label{sec:pilot}

Nine external raters labeled 20 blind A/B pairs (Fleiss'
$\kappa = 0.294$; 9 consensus pairs at $\geq$7/9 majority).
On these 9 pairs, harmonic compatibility---a signal with no
measurable weak-label contribution in our ablations---matches the
human majority in 7/9 cases, while tempo, the strongest
hard\_similar contributor in Table~\ref{tab:loso}, matches in
3/9---a descriptive reversal of the weak-label ranking.
The comparison remains exploratory: the pilot is post-hoc, small,
and uses a different corpus and pair construction than the weak-label
evaluation.  It raises the possibility of misalignment between proxy
rankings and human judgment but does not confirm it.

% ==================================================================
\section{Conclusion}\label{sec:conclusion}
% ==================================================================

Pocket Producer demonstrates that session-conditioned retrieval can
be operationalized in a DAW while keeping creative judgment with the
musician.
The deployed weak-label fusion is an \emph{operational} cold-start
candidate generator, but under clean pack-only labels it offers no
robust advantage over frozen cosine, its signal contributions switch
sign across negative regimes, and its rankings may not reflect human
preference.
Exactly this unreliability motivates the design: every recommendation
is previewable and reversible, and the proxy never decides alone.
Feedback logging provides the data foundation for future
personalization, which is not yet trained or evaluated.

% --- Page 3: references only ---

\newpage

\noindent\textbf{Acknowledgements.}
The author thanks the Audiotool team for the Nexus SDK and the FSLD
creators~\cite{ramires2020fsld}.

\vspace{0.5em}
\noindent\textbf{AI Usage Statement.}
Claude Code (Anthropic) was used for engineering implementation,
experiment scripting, and documentation drafting.
All research claims, experimental design, and analysis are the sole
responsibility of the author.

\bibliography{references}

\end{document}
